\section{References}
\label{sec:dd_references}%

The following references were consulted and used during the development of this Design Document:

\begin{enumerate}
    \item \textbf{SE2 R\&DD Assignment:} Project assignment document.  
    \item \textbf{Previous year student project:} Available on the WeBeep page for reference and guidance.
    \item \textbf{RASD of the Best Bike Paths project:} Requirements Analysis and Specification Document produced for the same project, used as the source of functional and non-functional requirements.
    \item \textbf{General Data Protection Regulation (GDPR):} Regulation on data protection and privacy.
    \item \textbf{RabbitMQ Message Brokers documentation:} Used for better understanding and representation of the message queuing.
    \item \textbf{OpenWeatherMap API documentation:} Used as an example of a weather provider for trip enrichment.
    \item \textbf{Map provider documentation for Google Maps Platform & Mapbox:} Used as reference for map visualization and geospatial services integration.
\end{enumerate}

\section{Used tools}
\label{sec:dd_used_tools}%

The following tools and technologies were used while producing this document:

\begin{enumerate}
    \item \textbf{GitHub} --- Version control and collaborative repository management.
    \item \textbf{\LaTeX{}} --- Document preparation and formatting with Overleaf as editing environment.
    \item \textbf{Visual Studio Code} --- Code editor used for writing and editing \LaTeX{} files and other project documentation.
    \item \textbf{Draw.io / PlantUML} --- Creation of UML and architectural diagrams.
    \item \textbf{Google Documents:} Collaborative platform for brainstorming, shared notes, and draft document preparation.
\end{enumerate}

\section{Use of Generative AI tools}
\label{sec:dd_ai_tools}%

Generative AI tools were used as support tools for specific, well-defined tasks during the development of this document. The following section describes the tools used, the inputs provided, and the outputs obtained.

\subsection*{Tool used}
\textbf{ChatGPT} ---  Large Language Model chatbot used for clarification of concepts, text rewriting, and assistance in structuring content.

\subsection*{Inputs and outputs}
\begin{enumerate}
    \item \textbf{Input:} ``Please rewrite and review this text as if you were a software expert. [copy of the text to be reviewed]''
    
    \textbf{Output (summary):} Rewritten text suggestions which were reviewed and integrated selectively to preserve technical intent and consistency.

    \item \textbf{Input:} ``Please check the diagrams in the document attached and make sure that all function calls are in the components interface list.''
    
    \textbf{Output (summary):} Most of the list was actually correct, the AI told me some functions that weren't included and also made a few mistakes.

    \item \textbf{Input:} ``How can I resize these three images so they are shown together in the same page?. [snippet of LaTeX code]''
    
    \textbf{Output (summary):} ChatGPT offered me a couple solutions to my problem so I could output the images as desired.

    \item \textbf{Input:} ``I am going to submit the work and I need you to review the entire work in general looking for typos and different things that may have been overlooked.''
    
    \textbf{Output (summary):} The AI tool identified and corrected multiple spelling errors, grammatical mistakes, and inconsistencies throughout the document, including typos, incorrect word usage, and duplicate labels.
\end{enumerate}

\textit{Note: Additional queries were used during the document preparation, but as they were redundant with the queries presented above, they have not been included in this document.}

\subsection*{Verification and refinement process}

All outputs from the generative AI tool were subject to the following verification process:
\begin{itemize}
    \item \textbf{Technical Accuracy:} All suggestions were cross-referenced with course materials, IEEE standards, and the project assignment specifications.
    \item \textbf{Content Validation:} Generated content was reviewed to ensure it accurately reflected the Best Bike Paths system requirements and did not introduce incorrect or misleading information.
    \item \textbf{Integration Review:} Before integration, all AI-generated suggestions were evaluated by both team members to ensure consistency with the document's overall structure and style.
    \item \textbf{Final Verification:} All integrated content was reviewed in the context of the complete document to ensure coherence and accuracy.
\end{itemize}

\subsection*{Responsibility statement}
The authors take full responsibility for all content in this document. The generative AI tool was used solely as a support tool for clarification, text improvement, and structural guidance. All final content, including text, diagrams, and specifications, was created, verified, and approved by the authors, who are fully capable of explaining and defending every aspect of the work presented in this document.